% GENERATED FILE. Edit canonical lessons, then run npm run generate:books.
% canonical-source-hash: fnv1a64:730fe5be86cd852f
% canonical-lessons: ES-C694-gaviota, ES-C694-ave, ES-C694-insecto, ES-C694-mariposa, ES-C694-lagartija

\chapter{Animales del campo}
\label{ch:es-animales-del-campo}
\hlchaptermodality{694}

\begin{chapteropening}
\textbf{By the end of this chapter:} \emph{I can name small wild creatures: seagulls, birds, insects, butterflies and lizards.}

Complete the last lesson of chapter 694: I can name small wild creatures: seagulls, birds, insects, butterflies and lizards.
\end{chapteropening}

\section[la gaviota]{la gaviota — the seagull}

\label{lesson:ES-C694-gaviota}



\begin{quote}
Before the new one: say the Spanish for the vegetation, then the Spanish for the oak tree.
\end{quote}

\subsection*{You'll want to know: la gaviota}
\textbf{la gaviota} — \textquotedblleft{}the seagull\textquotedblright{}.

\begin{quote}
Las gaviotas vuelan sobre el puerto.
\end{quote}

\begin{quote}
\emph{The seagulls fly over the harbour.}
\end{quote}

Common in ports and on beaches.

\begin{cousinweb}
\emph{gaviota} is built on Latin \emph{gāvia} \textquotedblleft{}a sea bird, perhaps the seamew\textquotedblright{} (Lewis \& Short; EtymDB).
\end{cousinweb}

\begin{grammarlens}[title={What you've built}]
One more word for telling, explaining and describing at B1.
\end{grammarlens}

\subsection*{Your turn}
\begin{itemize}
\raggedright
  \item \emph{Say it:} \emph{la gaviota}
  \item \emph{Say it:} \emph{la gaviota}, once more
  \item \emph{Say it:} a sentence of your own with \emph{la gaviota}
  \item \emph{Recall:} say the Spanish for the vegetation, then the Spanish for the oak tree, then say \emph{la gaviota} again
\end{itemize}

\begin{tcolorbox}[breakable,colback=teal!4,colframe=teal!35!black,title={Before you move on}]
What does la gaviota mean? (\textquotedblleft{}The seagull\textquotedblright{}.) Say it once more.
\end{tcolorbox}
\section[el ave]{el ave — the large bird, the bird (formal); poultry}

\label{lesson:ES-C694-ave}



\begin{quote}
Before the new one: say the Spanish for the oak tree, then the Spanish for the seagull.
\end{quote}

\subsection*{You'll want to know: el ave}
\textbf{el ave} — \textquotedblleft{}the large bird, the bird (formal); poultry\textquotedblright{}.

\begin{quote}
En el lago hay muchas aves.
\end{quote}

\begin{quote}
\emph{There are a lot of birds on the lake.}
\end{quote}

Feminine, but \emph{el ave} because of the stressed a; plural \emph{las aves}.

\begin{cousinweb}
\emph{ave} is inherited from Latin \emph{avis} \textquotedblleft{}a bird\textquotedblright{} (EtymDB).
\end{cousinweb}

\begin{grammarlens}[title={What you've built}]
One more word for telling, explaining and describing at B1.
\end{grammarlens}

\subsection*{Your turn}
\begin{itemize}
\raggedright
  \item \emph{Say it:} \emph{el ave}
  \item \emph{Say it:} \emph{el ave}, once more
  \item \emph{Say it:} a sentence of your own with \emph{el ave}
  \item \emph{Recall:} say the Spanish for the oak tree, then the Spanish for the seagull, then say \emph{el ave} again
\end{itemize}

\begin{tcolorbox}[breakable,colback=teal!4,colframe=teal!35!black,title={Before you move on}]
What does el ave mean? (\textquotedblleft{}The large bird, the bird (formal); poultry\textquotedblright{}.) Say it once more.
\end{tcolorbox}
\section[el insecto]{el insecto — the insect}

\label{lesson:ES-C694-insecto}



\begin{quote}
Before the new one: say the Spanish for the seagull, then the Spanish for the large bird.
\end{quote}

\subsection*{You'll want to know: el insecto}
\textbf{el insecto} — \textquotedblleft{}the insect\textquotedblright{}.

\begin{quote}
Por la noche hay muchos insectos.
\end{quote}

\begin{quote}
\emph{There are a lot of insects at night.}
\end{quote}

\emph{Una picadura de insecto} — an insect bite.

\begin{cousinweb}
\emph{insecto} is a learned borrowing from Latin \emph{insectum} \textquotedblleft{}cut into\textquotedblright{} (EtymDB), from \emph{in} and \emph{secāre} \textquotedblleft{}to cut\textquotedblright{}: an insect body looks cut into parts.
\end{cousinweb}

\begin{grammarlens}[title={What you've built}]
One more word for telling, explaining and describing at B1.
\end{grammarlens}

\subsection*{Your turn}
\begin{itemize}
\raggedright
  \item \emph{Say it:} \emph{el insecto}
  \item \emph{Say it:} \emph{el insecto}, once more
  \item \emph{Say it:} a sentence of your own with \emph{el insecto}
  \item \emph{Recall:} say the Spanish for the seagull, then the Spanish for the large bird, then say \emph{el insecto} again
\end{itemize}

\begin{tcolorbox}[breakable,colback=teal!4,colframe=teal!35!black,title={Before you move on}]
What does el insecto mean? (\textquotedblleft{}The insect\textquotedblright{}.) Say it once more.
\end{tcolorbox}
\section[la mariposa]{la mariposa — the butterfly}

\label{lesson:ES-C694-mariposa}



\begin{quote}
Before the new one: say the Spanish for the large bird, then the Spanish for the insect.
\end{quote}

\subsection*{You'll want to know: la mariposa}
\textbf{la mariposa} — \textquotedblleft{}the butterfly\textquotedblright{}.

\begin{quote}
En primavera hay muchas mariposas.
\end{quote}

\begin{quote}
\emph{There are a lot of butterflies in spring.}
\end{quote}

A moth is \emph{la polilla}.

\begin{cousinweb}
The origin of \emph{mariposa} is not settled; it is often explained as \emph{Mari, posa} \textquotedblleft{}Mary, settle down\textquotedblright{}, from a children’s song.
\end{cousinweb}

\begin{grammarlens}[title={What you've built}]
One more word for telling, explaining and describing at B1.
\end{grammarlens}

\subsection*{Your turn}
\begin{itemize}
\raggedright
  \item \emph{Say it:} \emph{la mariposa}
  \item \emph{Say it:} \emph{la mariposa}, once more
  \item \emph{Say it:} a sentence of your own with \emph{la mariposa}
  \item \emph{Recall:} say the Spanish for the large bird, then the Spanish for the insect, then say \emph{la mariposa} again
\end{itemize}

\begin{tcolorbox}[breakable,colback=teal!4,colframe=teal!35!black,title={Before you move on}]
What does la mariposa mean? (\textquotedblleft{}The butterfly\textquotedblright{}.) Say it once more.
\end{tcolorbox}
\section[la lagartija]{la lagartija — the small lizard, the wall lizard}

\label{lesson:ES-C694-lagartija}



\begin{quote}
Before the new one: say the Spanish for the insect, then the Spanish for the butterfly.
\end{quote}

\subsection*{You'll want to know: la lagartija}
\textbf{la lagartija} — \textquotedblleft{}the small lizard, the wall lizard\textquotedblright{}.

\begin{quote}
Una lagartija toma el sol en la pared.
\end{quote}

\begin{quote}
\emph{A lizard is sunning itself on the wall.}
\end{quote}

A bigger one is \emph{un lagarto}.

\begin{cousinweb}
\emph{lagartija} is built on \emph{lagarto} \textquotedblleft{}lizard\textquotedblright{}, which is inherited from Latin \emph{lacertus} \textquotedblleft{}a lizard\textquotedblright{} (Lewis \& Short; EtymDB).
\end{cousinweb}

\begin{grammarlens}[title={What you've built}]
One more word for telling, explaining and describing at B1.
\end{grammarlens}

\subsection*{Your turn}
\begin{itemize}
\raggedright
  \item \emph{Say it:} \emph{la lagartija}
  \item \emph{Say it:} \emph{la lagartija}, once more
  \item \emph{Say it:} a sentence of your own with \emph{la lagartija}
  \item \emph{Recall:} say the Spanish for the insect, then the Spanish for the butterfly, then say \emph{la lagartija} again
\end{itemize}

\begin{tcolorbox}[breakable,colback=teal!4,colframe=teal!35!black,title={Before you move on}]
What does la lagartija mean? (\textquotedblleft{}The small lizard, the wall lizard\textquotedblright{}.) Say it once more.
\end{tcolorbox}
